% BT589 -- Homology Separation LaTeX Insert
% Standalone manuscript-ready insert for the W33 preprint.

\subsection{Levi Homology versus Phase-Cover Fiber Homology}\label{subsec:levi-vs-fiber-homology}

There are two natural homological objects in the phase-cover construction, and
it is important not to identify them.  The protected W33 homology comes from the
point-line Levi graph of $W(3,3)$, not from the scalar fiber of the phase cover.

Let $L$ denote the W33 point-line Levi graph.  It has
\[
|V(L)|=80,\qquad |E(L)|=160,
\]
and hence
\[
\boxed{\beta_1(L)=160-80+1=81.}
\]
This is the protected cycle rank:
\[
\boxed{H_1=81.}
\]
The same Levi graph has $1620$ simple $8$-cycles, and the flag-edge incidence
count satisfies
\[
\boxed{160\cdot81=1620\cdot8=12960.}
\]
Thus $12960$ is a Levi support-incidence count derived from the $H_1=81$ cycle
geometry.

The scalar phase cover begins from this $12960$-element support base.  Each base
incidence has four nonzero ternary scalar lifts.  As a minimal fiber graph, those
four lifts form a square, so for one fiber
\[
V_{\rm fiber}=4,\qquad E_{\rm fiber}=4,\qquad \beta_1^{\rm fiber}=1.
\]
Across all base incidences this gives
\[
V=12960\cdot4=51840,\qquad E=12960\cdot4=51840,
\]
with $12960$ connected fiber components.  Therefore
\[
\boxed{\beta_1^{\rm fiber}=51840-51840+12960=12960.}
\]
Consequently
\[
\boxed{81\ne12960.}
\]
The correct relationship is layering rather than equality:
\[
\boxed{
H_1(L)=81
\quad\leadsto\quad
12960\text{ Levi support incidences}
\quad\leadsto\quad
51840\text{ phase-cover lifts}.
}
\]
The scalar cover doubles the nonzero phase sheets, but it does not replace the
Levi cycle homology.  The former is a fiber square over each support incidence;
the latter is the genuine incidence homology of the W33 point-line geometry.
